\section{Limits and what the test cannot show} \label{sec:limits}

\begin{itemize}
\item Whatever the outcome, Sector 110 cannot show whether A or B is a planet on TIC 4206066 rather than an eclipsing
  companion star bound to it and too close to resolve, or a planet transiting such a companion; that needs
  high-resolution imaging and radial velocities. It cannot give
  the mass of either object, which needs radial velocities, or a position better than a few arcseconds.
\item The archival sectors were inspected, the signals were selected after inspection (Section~\ref{sec:signals}), and
  one choice in the ephemeris of A was made after looking at the data (Section~\ref{sec:ephem}).
\item Among the rules set after archival measurements (Appendix~\ref{app:history}), three matter most. The edge zones were introduced after the archival dry runs had
  shown artefacts at data gaps; the two widths tried gave the same result, and the adopted width follows from the
  filter. The 99.9th and 0.1st percentiles were chosen to bring the measured false-alarm level down. This makes detection
  stricter and widens the consistency band slightly. The empirical localisation $p$ replaced a Gaussian $p$ that
  called too many on-target signals ``off target''. The empirical $p$ is more lenient for moderate offsets and gives
  false ``off target'' labels at about the nominal rate. None of these was tuned
  further (Appendix~\ref{app:limits}).
\item The false-detection levels come from trials that share data, so their bounds are approximate and are not a
  guarantee for Sector 110; they are frequencies of the noise model, not probabilities that a detected signal is a
  planet.
\item Signal B will most likely give O4 (Section~\ref{sec:outcomes}); its expected depth rests on one archival transit,
  and its Sector 98 depth is about 1.7 times larger.
\item Edge zones cost data: in a Sector 110 light curve with two or three segments, a transit has roughly a 6 to 10\,\%
  chance to fall in an edge zone and not be tested.
\item The weak 1.76\,d feature seen in Sector 98 (about 130\,ppm) is not part of this test: a period measured in one
  sector cannot predict Sector 110 epochs. Any look at it in Sector 110 is exploratory and will be labelled so, and so
  is any search away from the predicted times.
\end{itemize}
The full list of known limits, stated before the data, is in Appendix~\ref{app:limits}.
